\section*{Chapter \ref{ch:b2:reaction-enthalpy} --- Enthalpies of Reaction}

\begin{solution}{exo:b2:reaction-enthalpy:1}
$\Delta_r H^\circ = 3(-393.51) + 4(-285.83) - (-104.7) - 5(0) =
\qty{-2219.2}{kJ/mol}$, negative: \omterm{def:b2:reaction-enthalpy:exothermic}{exothermic}. (The measured combustion
enthalpy is the same, \qty{-2219.2}{kJ/mol}.)
\end{solution}

\begin{solution}{exo:b2:reaction-enthalpy:2}
Methane: $890.3/16.04 = \qty{55.5}{kJ/g}$, so \qty{55.5}{MJ} per
kilogram. Dihydrogen: $285.83/2.016 = \qty{141.8}{kJ/g}$,
\qty{141.8}{MJ/kg}: about 2.6 times more per kilogram (but far less per
litre, the gas being light).
\end{solution}

\begin{solution}{exo:b2:reaction-enthalpy:3}
$\Delta\nu_{\text{gas}} = 6 - 7.5 = -1.5$, so $\Delta_r H = \Delta_r U +
\Delta\nu_{\text{gas}}RT = -3264 - 1.5 \times 8.314 \times 298.15 \times
10^{-3} = \qty{-3267.7}{kJ/mol}$, in agreement with the value computed from
formation enthalpies, \qty{-3267.6}{kJ/mol}.
\end{solution}

\begin{solution}{exo:b2:reaction-enthalpy:4}
The target is the first reaction minus twice the second:
$\Delta_r H^\circ = -393.5 - 2(-283.0) = \qty{+172.5}{kJ/mol}$. It is
\omterm{def:b2:reaction-enthalpy:exothermic}{endothermic}: the hot coke cools as it turns carbon dioxide into carbon
monoxide.
\end{solution}

\begin{solution}{exo:b2:reaction-enthalpy:5}
Bonds broken: C--H and Cl--Cl, $416 + 243 = \qty{659}{kJ/mol}$; bonds
formed: C--Cl and H--Cl, $350 + 432 = \qty{782}{kJ/mol}$: estimate
\qty{-123}{kJ/mol}. From formation enthalpies: $-83.7 - 92.3 + 74.9 =
\qty{-101.1}{kJ/mol}$. The \qty{22}{kJ/mol} gap comes from the mean C--Cl
value, an average over several molecules, and from the C--H bond broken in
methane, stronger than the mean of the four.
\end{solution}

\begin{solution}{exo:b2:reaction-enthalpy:6}
At \qty{298}{K}: $-241.83 + 285.83 = \qty{44.00}{kJ/mol}$. Kirchhoff with
$\Delta C_p = 33.59 - 75.35 = \qty{-41.76}{J/(K.mol)}$:
$44.00 - 0.04176 \times 101.85 = \qty{39.75}{kJ/mol}$ at \qty{400}{K}. The
tables give $-242.85 + 282.59 = \qty{39.74}{kJ/mol}$: the constant heat
capacities are excellent over this interval.
\end{solution}

\begin{solution}{exo:b2:reaction-enthalpy:7}
Ionisation: $4.341 \times 96.485 = \qty{418.8}{kJ/mol}$; attachment:
$-3.613 \times 96.485 = \qty{-348.6}{kJ/mol}$. \omterm{def:b2:reaction-enthalpy:lattice-enthalpy}{Lattice enthalpy} $= 436.7 +
89.0 + 121.3 + 418.8 - 348.6 = \qty{717}{kJ/mol}$, less than the
\qty{787}{kJ/mol} of sodium chloride: \ce{K+} is larger than \ce{Na+}, the
ions are farther apart and attract each other less.
\end{solution}

\begin{solution}{exo:b2:reaction-enthalpy:8}
Heat capacity of the calorimeter: $C = 2.00/0.418 = \qty{4.785}{kJ/K}$.
Heat released: $4.785 \times 0.583 = \qty{2.789}{kJ}$, for $\xi =
\qty{0.0500}{mol}$ (the acid is limiting, the base in slight excess):
$\Delta_r H = -2.789/0.0500 = \qty{-55.8}{kJ/mol}$.
\end{solution}

\begin{solution}{exo:b2:reaction-enthalpy:9}
\ce{C8H18(l) + 25/2O2(g) -> 8CO2(g) + 9H2O(l)}: $\Delta_r H^\circ =
8(-393.51) + 9(-285.83) + 250.3 = \qty{-5470.3}{kJ/mol}$; with $M =
\qty{114.2}{g/mol}$, \qty{47.9}{MJ/kg}. Carbon dioxide: $8 \times 44.01 =
\qty{352.1}{g}$ for \qty{5.470}{MJ}, that is \qty{64.4}{g/MJ}, against
$44.01/0.8903 = \qty{49.4}{g/MJ}$ for methane: methane has more hydrogen per
carbon, and burning hydrogen releases heat without carbon dioxide.
\end{solution}

\begin{solution}{exo:b2:reaction-enthalpy:10}
Products of \ce{CH4 + 2O2 -> CO2 + 2H2O(g)}: $C_p = 37.13 + 2 \times 33.59 =
\qty{104.3}{J/K}$; $q = \qty{802.3}{kJ}$; $T_f = 298 + 802\,300/104.3 \approx
\qty{7990}{K}$, almost three times the value in air, because no nitrogen
takes its share of the heat. The real oxygen flame is far colder: long before \qty{7990}{K}, carbon dioxide and water
dissociate into carbon monoxide, hydrogen, oxygen and atoms, \omterm{def:b2:reaction-enthalpy:exothermic}{endothermic}
reactions that absorb most of the heat; and the heat capacities grow with
$T$.
\end{solution}

\begin{solution}{exo:b2:reaction-enthalpy:11}
$\Delta_r H^\circ(700) = \Delta_r H^\circ(298) + \int_{298}^{700}(\Delta_r a +
\Delta_r b\,T)\,\dd T = -91.88 + [-60.5 \times 401.85 + 0.0270 \times
(700^2 - 298.15^2)] \times 10^{-3} = -91.88 - 24.31 + 10.83 =
\qty{-105.4}{kJ/mol}$, within \qty{0.2}{kJ/mol} of the tables
(\qty{-105.2}{kJ/mol}), where the constant-$C_p$ estimate was off by
\qty{4.5}{kJ/mol}.
\end{solution}

\begin{solution}{exo:b2:reaction-enthalpy:12}
$\Delta_r H^\circ = 2 \times 90.29 = \qty{+180.6}{kJ/mol}$ (\omterm{def:b2:reaction-enthalpy:exothermic}{endothermic}).
$\Delta_r C_p^\circ = 2(29.85) - 29.12 - 29.38 = \qty{1.2}{J/(K.mol)}$, so
between \qty{298}{K} and \qty{2000}{K} the \omterm{def:b2:reaction-enthalpy:reaction-quantity}{reaction enthalpy} changes by
about $1.2 \times 1702 = \qty{2}{kJ/mol}$, one per cent. An \omterm{def:b2:reaction-enthalpy:exothermic}{endothermic}
reaction is favoured by high temperature (\cref{ch:b2:equilibrium-shifts}):
cold air makes almost no nitrogen monoxide, the hottest flames make the
most, which is why combustion temperature is a lever against this
pollutant.
\end{solution}

\begin{solution}{pb:b2:reaction-enthalpy:1}
\textbf{1.} \ce{C4H10(g) + 13/2O2(g) -> 4CO2(g) + 5H2O(l)}.
\textbf{2.} $\Delta_r H^\circ = 4(-393.51) + 5(-285.83) + 125.6 =
\qty{-2877.6}{kJ/mol}$.
\textbf{3.} With water vapour, $4(-393.51) + 5(-241.83) + 125.6 =
\qty{-2657.6}{kJ/mol}$; the difference, \qty{220.0}{kJ}, is the enthalpy of
vaporisation of five moles of water (\qty{44.0}{kJ/mol} each).
\textbf{4.} $2877.6/58.12 = \qty{49.5}{kJ/g}$ and $2657.6/58.12 =
\qty{45.7}{kJ/g}$.
\textbf{5.} The second: the water leaves as vapour and its heat of
condensation is lost.
\textbf{6.} $230 \times 45.7 = \qty{1.05e4}{kJ}$, about \qty{10.5}{MJ}.
\textbf{7.} At constant volume the heat received is $\Delta U$ (no
pressure work): the bomb measures $\Delta_r U$.
\textbf{8.} $\Delta\nu_{\text{gas}} = 4 - 1 - 6.5 = -3.5$ (the water is
liquid).
\textbf{9.} $\Delta_r U = \Delta_r H - \Delta\nu_{\text{gas}}RT = -2877.6 +
3.5 \times 2.479 = \qty{-2868.9}{kJ/mol}$.
\textbf{10.} $\xi = 0.500/58.12 = \qty{8.60e-3}{mol}$; heat $8.60 \times
10^{-3} \times 2868.9 = \qty{24.68}{kJ}$; $\Delta T = 24.68/10.00 =
\qty{2.468}{K}$.
\textbf{11.} Heat $10.00 \times 2.461 = \qty{24.61}{kJ}$: $\Delta_r U =
-24.61/(8.603 \times 10^{-3}) = \qty{-2860.6}{kJ/mol}$ and $\Delta_r H =
-2860.6 - 8.7 = \qty{-2869.3}{kJ/mol}$, \qty{0.29}{\percent} below the
tables in absolute value: heat lost by the calorimeter or a slightly
incomplete combustion.
\textbf{12.} $n = 1000/18.02 = \qty{55.5}{mol}$; $Q = 55.5 \times 75.35
\times 80 = \qty{334.5}{kJ}$.
\textbf{13.} Heat released $16.0 \times 45.72 = \qty{731.5}{kJ}$; efficiency
$334.5/731.5 = 0.46$.
\textbf{14.} Into the hot combustion gases that escape round the pan, the
air heated by the flame and the pan's walls, and radiation.
\textbf{15.} $230 \times 45.72 \times 0.457/334.5 = \qty{14.4}{L}$.
\textbf{16.} \qty{6.5}{mol} of \ce{O2} come with $6.5 \times 78.08/20.95 =
\qty{24.23}{mol}$ of \ce{N2} and $6.5 \times 0.93/20.95 = \qty{0.289}{mol}$
of \ce{Ar}.
\textbf{17.} \qty{4}{mol} \ce{CO2}, \qty{5}{mol} \ce{H2O(g)},
\qty{24.23}{mol} \ce{N2}, \qty{0.289}{mol} \ce{Ar}.
\textbf{18.} The flame is adiabatic and isobaric, so $\Delta H = 0$. Along a
path made of the reaction at \qty{298}{K} ($\Delta H_1 =
\qty{-2657.6}{kJ}$) followed by the heating of these products, $\Delta H_2 =
-\Delta H_1 = \qty{+2657.6}{kJ}$: all the heat goes into the gases.
\textbf{19.} $C_p = 4(37.13) + 5(33.59) + 24.23(29.12) + 0.289(20.79) =
\qty{1028}{J/K}$.
\textbf{20.} $T_f = 298 + 2\,657\,600/1028 = \qty{2883}{K}$.
\textbf{21.} $T_f = 2300 + 100 \times (2657.6 - 2515.7)/(2655.7 - 2515.7) =
\qty{2401}{K}$.
\textbf{22.} The heat capacities of the gases grow with temperature (by
about a quarter between \qty{298}{K} and \qty{2400}{K}): the room-temperature
values heat the gases too cheaply.
\textbf{23.} Above \qty{2000}{K} carbon dioxide and water partly dissociate
(\omterm{def:b2:reaction-enthalpy:exothermic}{endothermic}), the flame radiates, and mixing with air is never perfect.
\textbf{24.} The \omterm{def:b2:reaction-enthalpy:flame-temperature}{adiabatic flame temperature} of butane in air is
$\boldsymbol{T_{\text{ad}} \approx \qty{2.40e3}{K}}$.
\end{solution}
